\documentclass[UTF8,a4paper,11pt]{ctexart}

\usepackage{geometry}
\geometry{left=2.2cm,right=2.2cm,top=2.4cm,bottom=2.4cm,headheight=14pt,headsep=0.4cm}

\usepackage{graphicx}
\usepackage{booktabs}
\usepackage{array}
\usepackage{tabularx}
\usepackage{longtable}
\usepackage{enumitem}
\usepackage{xcolor}
\usepackage{fancyhdr}
\usepackage{titlesec}
\usepackage{tcolorbox}
\usepackage{tikz}
\usetikzlibrary{arrows.meta,positioning,fit,calc}

\definecolor{main}{RGB}{38,70,83}
\definecolor{sub}{RGB}{233,236,239}

\pagestyle{fancy}
\fancyhf{}
\lhead{实验5 存储器及数据通路设计}
\rhead{\thepage}
\renewcommand{\headrulewidth}{0.4pt}

\titleformat{\section}{\Large\bfseries\color{main}}{\thesection}{0.6em}{}
\titleformat{\subsection}{\large\bfseries\color{main}}{\thesubsection}{0.5em}{}

\setlist[itemize]{leftmargin=2em,itemsep=0.3em,topsep=0.3em}
\setlist[enumerate]{leftmargin=2em,itemsep=0.3em,topsep=0.3em}

\newcommand{\blankline}[1]{\vspace{0.25em}\noindent\rule{#1}{0.4pt}\vspace{0.25em}}
\newcommand{\placeholder}[1]{
\begin{tcolorbox}[colback=sub,colframe=main,boxrule=0.6pt,arc=2mm,left=2mm,right=2mm,top=1mm,bottom=1mm]
#1
\end{tcolorbox}
}
\newcommand{\truthtable}[4]{%
\par\medskip
\noindent\begin{tcolorbox}[colback=white,colframe=main,boxrule=0.8pt,arc=2mm,left=1.5mm,right=1.5mm,top=1mm,bottom=1mm]
\centering\bfseries #1\par\smallskip
\renewcommand{\arraystretch}{0.75}
\begin{tabular}{#2}
\toprule
#3\\
\midrule
#4
\bottomrule
\end{tabular}
\end{tcolorbox}
\par\medskip%
}
\newcommand{\figureplaceholder}[2]{
\par\medskip
\noindent\begin{tcolorbox}[colback=white,colframe=main,boxrule=0.8pt,arc=2mm,left=1.5mm,right=1.5mm,top=1mm,bottom=1mm]
\centering\bfseries #1\par\smallskip
\IfFileExists{figures/#2}{%
\includegraphics[width=0.80\linewidth,height=0.88\textheight,keepaspectratio]{figures/\detokenize{#2}}%
}{%
\fbox{\parbox[c][6cm][c]{0.80\linewidth}{\centering 图片请放在 figures/\texttt{\detokenize{#2}}}}
}%
\end{tcolorbox}
\par\medskip
}
\newcommand{\riskimage}[2]{%
\begin{minipage}[t]{0.31\linewidth}
\centering
\includegraphics[height=3.75cm]{figures/\detokenize{#1}}\par
{\small #2}
\end{minipage}%
}

\begin{document}

\begin{center}
{\Huge\bfseries 实验报告}\\[0.8em]
{\LARGE 实验5 存储器及数据通路设计}\\[1.2em]
\end{center}

\begin{tcolorbox}[colback=white,colframe=main,boxrule=0.8pt,arc=2mm]
\renewcommand{\arraystretch}{1.6}
\begin{tabularx}{\textwidth}{>{\bfseries}p{2.6cm}X >{\bfseries}p{2.6cm}X}
课程名称 & 数字逻辑与计算机组成 & 指导教师 & 吴海军 \\
\end{tabularx}
\end{tcolorbox}

\section{实验目的}
\placeholder{\begin{itemize}
\item 掌握利用 ROM 组件实现指令存储器的设计方法。
\item 掌握 ASCII 点阵字库的使用方法。
\item 掌握利用 RAM 组件实现按字节存取的数据存储器的设计方法。
\item 掌握实现 RV32I 存取指令的设计方法。
\item 掌握 RV32I 下地址逻辑部件的设计方法。
\item 掌握 RV32I 立即数扩展器的设计方式。
\item 根据 RV32I 指令执行过程，掌握 RV32I 运算类指令数据通路的设计方法。
\item 掌握跳转控制器的设计方法。
\end{itemize}}

\section{实验环境}
\placeholder{\begin{itemize}
\item 平台：Logisim。
\end{itemize}}

\section{实验内容}

\subsection{只读存储器实验}

\subsubsection{实验整体方案设计}
\placeholder{\begin{itemize}
\item 顶层模块设计：实验电路较为简单，不需要顶层模块设计图。
\item 输入输出引脚：
\end{itemize}}

\begin{center}
\renewcommand{\arraystretch}{1.25}
\setlength{\tabcolsep}{8pt}
\begin{tabularx}{0.96\linewidth}{|>{\centering\arraybackslash}p{0.28\linewidth}|>{\centering\arraybackslash}X|}
\hline
ASCII-CODE & 输入显示字符的 ASCII 码 \\
\hline
row1 row15 & 输出信号 \\
\hline
CodeErr & 若 code 不在可见字符范围，则 CodeErr 为 1，不显示字符 \\
\hline
\end{tabularx}
\end{center}

\subsubsection{实验原理图和电路图}
\figureplaceholder{只读存储器电路图}{exp5_rom_circuit.png}

\subsubsection{实验数据仿真测试图}
\figureplaceholder{只读存储器仿真测试图}{exp5_rom_simulation.png}

\subsubsection{错误现象及分析}
\placeholder{
    本次实验中没有出现错误。
}

\subsection{数据存储器实验}

\subsubsection{实验整体方案设计}
\placeholder{\begin{itemize}
\item 顶层模块设计：实验电路较为简单，不需要顶层模块设计图。
\item 输入输出引脚：见下第一张表。
\item 关键控制信号与地址对齐关系：见下第二张表。
\end{itemize}}

\begin{center}
\renewcommand{\arraystretch}{1.25}
\setlength{\tabcolsep}{8pt}
\begin{tabularx}{0.96\linewidth}{|>{\centering\arraybackslash}p{0.28\linewidth}|>{\centering\arraybackslash}X|}
\hline
Addr & 地址 \\
\hline
CLK & 时钟信号 \\
\hline
WE & 写使能信号 \\
\hline
CLR & 清零信号 \\
\hline
Din & 输入数据 \\
\hline
MemOp & 控制信号 \\
\hline
Dout & 输出数据 \\
\hline
AlignErr & 错误信号（地址未对齐时为 1） \\
\hline
\end{tabularx}
\end{center}

\begin{center}
\renewcommand{\arraystretch}{1.25}
\setlength{\tabcolsep}{8pt}
\begin{tabularx}{0.98\linewidth}{|>{\centering\arraybackslash}p{0.14\linewidth}|>{\centering\arraybackslash}p{0.13\linewidth}|X|}
\hline
MemOp & 指令 & 含 义 \\
\hline
000 & lw,sw & 按字存取，4 字节 \\
\hline
001 & lbu & 按字节读取，1 个字节，0 扩展到 4 字节 \\
\hline
010 & lhu & 按半字读取，2 字节，0 扩展到 4 字节 \\
\hline
101 & lb,sb & 按字节存取，1 个字节，在读取时，按符号位扩展到 4 字节 \\
\hline
110 & lh,sh & 按半字存取，2 字节，在读取时，按符号位扩展到 4 字节 \\
\hline
\end{tabularx}
\end{center}

\subsubsection{实验原理图和电路图}
\figureplaceholder{数据存储器电路图}{exp5_dataram_circuit.png}

\subsubsection{实验数据仿真测试图}
\figureplaceholder{数据存储器仿真测试图}{exp5_dataram_simulation.png}

\subsubsection{错误现象及分析}
\placeholder{
    实验中没有出现错误。
}

\subsection{取指令部件 IFU 实验}

\subsubsection{实验整体方案设计}
\placeholder{\begin{itemize}
\item 顶层模块设计：实验电路较为简单，不需要顶层模块设计图。
\item 输入输出引脚：见下第一张表。
\item 关键控制信号：见下第二张表。
\end{itemize}}

\begin{center}
\renewcommand{\arraystretch}{1.25}
\setlength{\tabcolsep}{8pt}
\begin{tabularx}{0.96\linewidth}{|>{\centering\arraybackslash}p{0.28\linewidth}|>{\centering\arraybackslash}X|}
\hline
PC & 程序计数器 \\
\hline
Initial Address & 初始地址 \\
\hline
Reset & 复位初始地址 \\
\hline
Imm & 程序跳转的偏移量（立即数） \\
\hline
BusA & 寄存器 Rs1 中的数据 \\
\hline
CLK & 时钟信号 \\
\hline
NPCASrc & 控制 PC 加法器输入端 A 的信号 \\
\hline
NPCBSrc & 控制 PC 加法器输入端 B 的信号 \\
\hline
\end{tabularx}
\end{center}

\begin{center}
\renewcommand{\arraystretch}{1.25}
\setlength{\tabcolsep}{7pt}
\begin{tabularx}{0.98\linewidth}{|>{\centering\arraybackslash}p{0.18\linewidth}|>{\centering\arraybackslash}p{0.12\linewidth}|>{\centering\arraybackslash}p{0.12\linewidth}|X|}
\hline
指令类型 & NPCASrc & NPCBSrc & 下条指令地址 \\
\hline
顺序指令 & 0 & 0 & PC=PC+4 \\
\hline
无条件跳转 jal & 0 & 1 & PC=PC+Imm \\
\hline
无条件跳转 jalr & 1 & 1 & PC=R[rs1]+Imm \\
\hline
分支跳转指令 & 0 & 1 & 条件成立 PC=PC+Imm，否则 PC=PC+4 \\
\hline
\end{tabularx}
\end{center}

\subsubsection{实验原理图和电路图}
\figureplaceholder{IFU 取指令部件电路图}{exp5_ifu_circuit.png}

\subsubsection{实验数据仿真测试图}
\figureplaceholder{IFU 取指令部件仿真测试图}{exp5_ifu_simulation.png}

\begingroup
\renewcommand{\arraystretch}{1.25}
\setlength{\tabcolsep}{0.5pt}
\setlength{\LTleft}{0pt}
\setlength{\LTright}{0pt}
\begin{longtable}{|>{\centering\arraybackslash}p{0.8cm}|>{\centering\arraybackslash}p{1.0cm}|>{\centering\arraybackslash}p{1.5cm}|>{\centering\arraybackslash}p{2.0cm}|>{\centering\arraybackslash}p{1.45cm}|>{\centering\arraybackslash}p{1.7cm}|>{\centering\arraybackslash}p{1.7cm}|>{\centering\arraybackslash}p{1.6cm}|>{\centering\arraybackslash}p{4.1cm}|}
\hline
\multicolumn{1}{|c|}{} & \multicolumn{6}{c|}{输入} & \multicolumn{2}{c|}{输出（边沿信号有效前）} \\
\hline
序号 & Reset & Initial Address & Imm & BusA & NPCASrc & NPCBSrc & PC & IR \\
\hline
\endfirsthead
\hline
\endhead
\hline
\endfoot
\hline
\endlastfoot
1 & 1 & 0x400 & 0 & 0 & 0 & 0 & 0 & 0 \\
\hline
2 & 0 & 0 & 0 & 0 & 0 & 0 & 0x400 & 0x2083 \\
\hline
3 & 0 & 0 & 0 & 0 & 0 & 0 & 0x404 & 0x8133 \\
\hline
4 & 0 & 0 & 0x7f8 & 0 & 0 & 1 & 0x408 & 0x106213 \\
\hline
5 & 0 & 0 & 0x20c & 0 & 0 & 1 & 0xc00 & 0x413 \\
\hline
6 & 1 & 0x1500 & 0 & 0 & 0 & 0 & 0xe0c & 0xdeadbeef \\
\hline
7 & 0 & 0 & 0 & 0 & 0 & 0 & 0x1500 & 0x413 \\
\hline
8 & 0 & 0 & 0 & 0 & 0 & 0 & 0x1504 & 0x9117 \\
\hline
9 & 0 & 0 & 0 & 0x3fbc0 & 1 & 1 & 0x1508 & 0xffc10113 \\
\hline
10 & 0 & 0 & 0 & 0 & 0 & 0 & 0x3fbc0 & 0x413 \\
\hline
11 & 0 & 0 & 0x1fc & 0 & 0 & 1 & 0x3fbc4 & 0x9117 \\
\hline
12 & 0 & 0 & 0 & 0x150c & 1 & 1 & 0x3fdc0 & 0x40513 \\
\hline
13 & 0 & 0 & 0x0c00 & 0 & 1 & 1 & 0x150c & 0x39c000ef \\
\hline
14 & 0 & 0 & 0 & 0 & 0 & 0 & 0xc00 & 0x413 \\
\hline
15 & 0 & 0 & 0 & 0 & 0 & 0 & 0xc04 & 0x9117 \\
\hline
16 & 0 & 0 & 0xfffff9c8 & 0 & 0 & 1 & 0xc08 & 0xffc10113 \\
\hline
17 &  &  &  &  &  &  & 0x5d0 & 0xdeadbeef \\
\hline
\end{longtable}
\endgroup

\subsubsection{错误现象及分析}
\placeholder{
    本实验中没有出现错误。
}

\subsection{取操作数部件 IDU 实验}

\subsubsection{实验整体方案设计}
\placeholder{\begin{itemize}
\item 顶层模块设计：本实验电路较为简单，不需要顶层模块设计图。
\item 输入输出引脚：见下第一张表。
\end{itemize}}

\begin{center}
\renewcommand{\arraystretch}{1.25}
\setlength{\tabcolsep}{8pt}
\begin{tabularx}{0.96\linewidth}{|>{\centering\arraybackslash}p{0.28\linewidth}|>{\centering\arraybackslash}X|}
\hline
IR & 指令 \\
\hline
BusW & 写入寄存器堆的数据 \\
\hline
PC & 程序计数器 \\
\hline
ExtOp & 立即数扩展器控制信号 \\
\hline
RegWr & 寄存器堆写使能控制信号 \\
\hline
ALUASrc & ALU 输入端 A 的来源 \\
\hline
ALUBSrc & ALU 输入端 B 的来源 \\
\hline
CLK & 时钟信号 \\
\hline
DataA & ALU 输入端 A 的来源 \\
\hline
DataB & ALU 输入端 B 的来源 \\
\hline
BusA & 寄存器 Rs1 中的数据 \\
\hline
BusB & 寄存器 Rs2 中的数据 \\
\hline
Imm & 立即数 \\
\hline
OpCode & 操作码 \\
\hline
Func3,Func7 & 功能码 \\
\hline
\end{tabularx}
\end{center}

\subsubsection{实验原理图和电路图}
\figureplaceholder{IDU 取操作数部件电路图}{exp5_idu_circuit.png}
\figureplaceholder{IDU 立即数扩展器电路图}{exp5_idu_immext.png}
\placeholder{
    以下为 IDU 寄存器堆在本实验中做出的修改部分，使得第 0 号寄存器始终为 0，且在写入时，如果写入地址为 0，则不进行写入操作。其余均与先前实验中设计的寄存器堆相同。
}
\figureplaceholder{IDU 寄存器堆修改部分电路图}{exp5_idu_regfile.png}

\subsubsection{实验数据仿真测试图}
\placeholder{
    以下为取操作数部件输入输出数据列表。
}

\begingroup
\small
\renewcommand{\arraystretch}{1.15}
\setlength{\tabcolsep}{1pt}
\setlength{\LTleft}{0pt}
\setlength{\LTright}{0pt}
\begin{longtable}{|>{\centering\arraybackslash}p{0.8cm}|>{\centering\arraybackslash}p{1.85cm}|>{\centering\arraybackslash}p{2.05cm}|>{\centering\arraybackslash}p{3.1cm}|>{\centering\arraybackslash}p{3.95cm}|>{\centering\arraybackslash}p{3.05cm}|}
\hline
PC & IR & BusW & 控制信号值 & 指令功能 / 汇编语句 & DataA、DataB 输出值 \\
\hline
\endfirsthead
\hline
PC & IR & BusW & 控制信号值 & 指令功能 / 汇编语句 & DataA、DataB 输出值 \\
\hline
\endhead
\hline
\endfoot
\hline
\endlastfoot
0x0 & 0xfedca2b7 & 0xfedca000 & ExtOp=1 RegWr=1 ALUASrc=0 ALUBSrc=2 & lui x5,0xfedca & 0x0 / 0xfedca000 \\
\hline
0x4 & 0xf9c28293 & 0xfedc9f9c & ExtOp=0 RegWr=1 ALUASrc=0 ALUBSrc=2 & addi x5,x5,-100 & 0xfedca000 / 0xffffff9c \\
\hline
0x8 & 0x01006013 & 0x10 & ExtOp=0 RegWr=1 ALUASrc=0 ALUBSrc=2 & ori x0,x0,0x10 & 0x0 / 0x10 \\
\hline
0xc & 0x06502223 & 0x64 & ExtOp=2 RegWr=0 ALUASrc=0 ALUBSrc=2 & sw x5,0x64(x0) & 0x0 / 0x64 \\
\hline
0x10 & 0x06400303 & 0xffffff9c & ExtOp=0 RegWr=1 ALUASrc=0 ALUBSrc=2 & lb x6,0x64(x0) & 0x0 / 0x64 \\
\hline
0x14 & 0x06601423 & 0x68 & ExtOp=2 RegWr=0 ALUASrc=0 ALUBSrc=2 & sh x6,0x68(x0) & 0x0 / 0x68 \\
\hline
0x18 & 0x06405383 & 0x9f9c & ExtOp=0 RegWr=1 ALUASrc=0 ALUBSrc=2 & lhu x7,0x64(x0) & 0x0 / 0x64 \\
\hline
0x1c & 0x06701623 & 0x6c & ExtOp=2 RegWr=0 ALUASrc=0 ALUBSrc=2 & sh x7,0x6c(x0) & 0x0 / 0x6c \\
\hline
0x20 & 0x4042d413 & 0xffedc9f9 & ExtOp=0 RegWr=1 ALUASrc=0 ALUBSrc=2 & srai x8,x5,0x4 & 0xfedc9f9c / 0x404 \\
\hline
0x24 & 0x006444b3 & 0x123665 & ExtOp=0 RegWr=1 ALUASrc=0 ALUBSrc=0 & xor x9,x8,x6 & 0xffedc9f9 / 0xffffff9c \\
\hline
0x28 & 0x00649533 & 0x50000000 & ExtOp=0 RegWr=1 ALUASrc=0 ALUBSrc=0 & sll x10,x9,x6 & 0x123665 / 0xffffff9c \\
\hline
0x2c & 0x008505b3 & 0x4fedc9f9 & ExtOp=0 RegWr=1 ALUASrc=0 ALUBSrc=0 & add x11,x10,x8 & 0x50000000 / 0xffedc9f9 \\
\hline
0x30 & 0x00b2a633 & 0x1 & ExtOp=0 RegWr=1 ALUASrc=0 ALUBSrc=0 & slt x12,x5,x11 & 0xfedc9f9c / 0x4fedc9f9 \\
\hline
0x34 & 0x00b2b6b3 & 0x0 & ExtOp=0 RegWr=1 ALUASrc=0 ALUBSrc=0 & sltu x13,x5,x11 & 0xfedc9f9c / 0x4fedc9f9 \\
\hline
0x38 & 0x40b287b3 & 0xaeeed5a3 & ExtOp=0 RegWr=1 ALUASrc=0 ALUBSrc=0 & sub x15,x5,x11 & 0xfedc9f9c / 0x4fedc9f9 \\
\hline
0x3c & 0x06f02823 & 0x70 & ExtOp=2 RegWr=0 ALUASrc=0 ALUBSrc=2 & sw x15,0x70(x0) & 0x0 / 0x70 \\
\hline
0x40 & 0x0067c263 & 0x1 & ExtOp=3 RegWr=0 ALUASrc=0 ALUBSrc=0 & blt x15,x6,0x4 & 0xaeeed5a3 / 0xffffff9c \\
\hline
0x44 & 0x0067d263 & 0x1 & ExtOp=3 RegWr=0 ALUASrc=0 ALUBSrc=0 & bge x15,x6,0x4 & 0xaeeed5a3 / 0xffffff9c \\
\hline
0x48 & 0x0040086f & 0x4c & ExtOp=4 RegWr=1 ALUASrc=1 ALUBSrc=1 & jal x16,0x4 & 0x48 / 0x4 \\
\hline
0x4c & 0x004808e7 & 0x50 & ExtOp=0 RegWr=1 ALUASrc=1 ALUBSrc=1 & jalr x17,x16,0x4 & 0x4c / 0x4 \\
\hline
0x50 & 0x00001917 & 0x1050 & ExtOp=1 RegWr=1 ALUASrc=1 ALUBSrc=2 & auipc x18,0x1 & 0x50 / 0x1000 \\
\hline
\end{longtable}
\endgroup
\figureplaceholder{取操作数部件 IDU 仿真测试图}{exp5_idu_simulation.png}

\subsubsection{错误现象及分析}
\placeholder{
    一开始的错误在于将两种立即数拓展方式的控制型号理解反了，使得前两组测试数据的输出与答案不符。在添加了一个非门后，问题得以解决。
}

\subsection{数据通路实验}

\subsubsection{实验整体方案设计}
\placeholder{\begin{itemize}
\item 顶层模块设计：本实验电路较为简单，不需要顶层模块设计图。
\item 主要子电路：IFU、IDU、ALU、DataRAM、BandJ、寄存器堆等。
\end{itemize}}

\subsubsection{实验原理图和电路图}
\figureplaceholder{单周期 CPU 数据通路电路图}{exp5_datapath_circuit.png}
\figureplaceholder{跳转控制器组件电路图}{exp5_datapath_bj.png}

\subsubsection{实验数据仿真测试图}
\figureplaceholder{数据通路仿真测试图}{exp5_datapath_simulation.png}

\begingroup
\scriptsize
\renewcommand{\arraystretch}{1.03}
\setlength{\tabcolsep}{1pt}
\setlength{\LTleft}{0pt}
\setlength{\LTright}{0pt}
\begin{longtable}{|>{\centering\arraybackslash}p{0.65cm}|>{\centering\arraybackslash}p{1.55cm}|>{\centering\arraybackslash}p{6.25cm}|>{\centering\arraybackslash}p{1.3cm}|>{\centering\arraybackslash}p{1.6cm}|>{\centering\arraybackslash}p{1.55cm}|>{\centering\arraybackslash}p{1.3cm}|>{\centering\arraybackslash}p{1.3cm}|}
\hline
序号 & 指令 & 控制信号赋值，未列出的控制信号值为0。 & \multicolumn{5}{c|}{输出信号值} \\
\cline{4-8}
& & & PC & BusW & Din & NPCASrc & NPCBSrc \\
\hline
\endfirsthead
\hline
序号 & 指令 & 控制信号赋值，未列出的控制信号值为0。 & \multicolumn{5}{c|}{输出信号值} \\
\cline{4-8}
& & & PC & BusW & Din & NPCASrc & NPCBSrc \\
\hline
\endhead
\hline
\endfoot
\hline
\endlastfoot
1 & fedca2b7 & ExtOp=1, RegWr=1, ALUBSrc=2, ALUctr=f & 0x400 & 0xfedca000 & 0x0 & 0 & 0 \\
\hline
2 & f9c28293 & RegWr=1, ALUBSrc=2 & 0x404 & 0xfedc9f9c & 0x0 & 0 & 0 \\
\hline
3 & 01006013 & RegWr=1, ALUBSrc=2, ALUctr=6 & 0x408 & 0x10 & 0x0 & 0 & 0 \\
\hline
4 & 06502223 & ExtOp=2, MemWr=1, ALUBSrc=2 & 0x40c & 0x64 & 0xfedc9f9c & 0 & 0 \\
\hline
5 & 06400303 & MemOp=5, RegWr=1, ALUBSrc=2, MemtoReg=1 & 0x410 & 0xffffff9c & 0x0 & 0 & 0 \\
\hline
6 & 06601423 & ExtOp=2, MemWr=1, MemOp=6, ALUBSrc=2 & 0x414 & 0x68 & 0xffffff9c & 0 & 0 \\
\hline
7 & 06405383 & MemOp=2, RegWr=1, ALUBSrc=2, MemtoReg=1 & 0x418 & 0x9f9c & 0x0 & 0 & 0 \\
\hline
8 & 06701623 & ExtOp=2, MemWr=1, MemOp=6, ALUBSrc=2 & 0x41c & 0x6c & 0x9f9c & 0 & 0 \\
\hline
9 & 4042d413 & RegWr=1, ALUBSrc=2, ALUctr=d & 0x420 & 0xffedc9f9 & 0x0 & 0 & 0 \\
\hline
10 & 006444b3 & RegWr=1, ALUctr=4 & 0x424 & 0x123665 & 0xffffff9c & 0 & 0 \\
\hline
11 & 00649533 & RegWr=1, ALUctr=1 & 0x428 & 0x50000000 & 0xffffff9c & 0 & 0 \\
\hline
12 & 008505b3 & RegWr=1 & 0x42c & 0x4fedc9f9 & 0xffedc9f9 & 0 & 0 \\
\hline
13 & 00b2a633 & RegWr=1, ALUctr=2 & 0x430 & 0x1 & 0x4fedc9f9 & 0 & 0 \\
\hline
14 & 00b2b6b3 & RegWr=1, ALUctr=3 & 0x434 & 0x0 & 0x4fedc9f9 & 0 & 0 \\
\hline
15 & 40b287b3 & RegWr=1, ALUctr=8 & 0x438 & 0x0 & 0x4fedc9f9 & 0 & 0 \\
\hline
16 & 06f02823 & ExtOp=2, MemWr=1, ALUBSrc=2 & 0x43c & 0x70 & 0x0 & 0 & 0 \\
\hline
17 & 0067c263 & ExtOp=3, ALUctr=2, BandJ=6 & 0x440 & 0x0 & 0xffffff9c & 0 & 0 \\
\hline
18 & 0067d263 & ExtOp=3, ALUctr=2, BandJ=7 & 0x444 & 0x0 & 0xffffff9c & 0 & 1 \\
\hline
19 & 0040086f & ExtOp=4, RegWr=1, ALUABSrc=1, ALUBSrc=1, BandJ=1 & 0x448 & 0x44c & 0x0 & 0 & 1 \\
\hline
20 & 004808e7 & RegWr=1, ALUABSrc=1, ALUBSrc=1, BandJ=2 & 0x44c & 0x450 & 0x0 & 1 & 1 \\
\hline
21 & 00001917 & ExtOp=1, RegWr=1, ALUABSrc=1, ALUBSrc=2 & 0x450 & 0x1450 & 0x0 & 0 & 0 \\
\hline
\end{longtable}
\endgroup

\subsubsection{错误现象及分析}
\placeholder{
    错误来源于加载了错误的指令数据。在加载正确的指令数据后，正确通过了所有测试数据。
}

\section{思考题}

\subsection{如何利用 ROM 实验实现滚动显示的功能，在 3 个 LED 点阵矩阵中，左右滚动显示 5 个 ASCII 字符？}
\placeholder{
    可以在 ROM 中存储 5 个 ASCII 字符的点阵字模数据，并设计一个计数器来控制当前显示的字符位置。每当计数器增加时，显示的字符就向左滚动一个位置。当计数器达到一定值（0x4f）时，重新回到初始位置，实现循环滚动显示的效果。电路文件存储在 lab5.6.circ 中。
}
\figureplaceholder{滚动显示布局图}{exp5_scroll_design.png}

\subsection{分析说明如果 PC 寄存器、寄存器堆和数据存储器写入数据的时钟信号触发边沿不一致，则对程序执行有何影响？}
\placeholder{
    如果 PC、寄存器堆和数据存储器的写入边沿不一致，就会破坏单周期 CPU 中“一个周期内状态同时更新”的假设，使各部件的写入与下一条指令的取数、译码不同步。这样容易产生读到旧值或未稳定数据的情况，造成数据冒险、写回错误、访存错误，严重时还会出现跳转地址错误或程序执行顺序紊乱。因此，通常应让这些时序部件尽量使用同一触发边沿，或通过严格的两相时钟设计保证时序正确。
}

\subsection{在 CPU 启动执行后，如何实现在当前程序结束后，CPU 不再继续往下执行？}
\placeholder{
    可以在程序末尾加入一条结束指令，例如 ecall，并在控制器中把它译码为停机信号，使 PCWrite 失效，从而冻结 PC，不再继续取后续指令；同时也可关闭寄存器堆和数据存储器的写使能，避免产生无意义的状态更新。若电路中没有专门的停机控制，也可以把最后一条指令写成跳转到自身的指令，例如 jal x0, 0，让 CPU 进入无限循环，从效果上停止继续向下执行。
}

\end{document}
